\documentclass[11pt]{article}
\usepackage{amsmath,amssymb,amsthm}
\usepackage[margin=1in]{geometry}
\usepackage{booktabs}
\usepackage[colorlinks=true,linkcolor=blue,urlcolor=blue]{hyperref}

\newtheorem{definition}{Definition}
\newtheorem{remark}{Remark}

\newcommand{\R}{\mathbb{R}}
\newcommand{\sg}{\operatorname{sg}}
\newcommand{\softmax}{\operatorname{softmax}}
\newcommand{\1}{\mathbf{1}}
\DeclareMathOperator{\Cat}{Cat}

\title{\bf The Cue-Triggered Lick-Timing Task:\\ Environment, Agent and Learning Rule}
\author{\texttt{timingtask}}
\date{\today}

\begin{document}
\maketitle

\noindent
This document specifies the reinforcement-learning formulation implemented in
\texttt{timingtask}: the observation, the recurrent agent, the reward, and the
computation performed by one training update. The reference implementation is
\texttt{generator.py} (environment), \texttt{models.py} (network) and \texttt{rl.py}
(agent and trainer). Default values refer to \texttt{variants.BASE}.

\tableofcontents

\section{Notation}\label{sec:notation}

\begin{definition}[Indices and dimensions]
\begin{itemize}
  \item $b \in \{1,\dots,B\}$ indexes the environment copy. $B$ independent trial
        generators run in lockstep and share one set of network weights; each keeps its
        own trial history, outcome history and delay schedule. Equations are written for a
        single $b$ except where the batch is summed over.
  \item $t \in \{0,\dots,T-1\}$ indexes the time step within one episode. One episode is
        one trial. The step size is $\mathrm{d}t$ seconds; $T$ is the length of the longest
        trial in the batch and shorter trials are padded (\S\ref{sec:mask}).
  \item $i$ indexes the gradient update. One update consumes one batch of $B$ episodes.
  \item $N$ is the number of recurrent units, $d_u$ the observation dimension and
        $|\mathcal{A}| = 2$ the number of actions (wait, lick).
\end{itemize}
\end{definition}

\begin{definition}[Stop-gradient]
$\sg[\cdot]$ is the identity in the forward pass and the zero map in the backward pass
(\texttt{.detach()} in PyTorch).
\end{definition}

\section{Observation}\label{sec:obs}

\begin{definition}[Observation vector]
At step $t$ the agent receives $u_t \in \R^{d_u}$ consisting of one or two cue channels
followed by $4K$ previous-trial channels, where $K$ is \texttt{n\_lags}:
\begin{align}
  c_t        &= \1\bigl[\,0 \le t - t_{\mathrm{cue}} < n_{\mathrm{cue}}\,\bigr]
    && \text{transient cue (\texttt{cue\_mode} pulse or both)}, \\
  s_t        &= \1\bigl[\, t \ge t_{\mathrm{cue}}\,\bigr]
    && \text{tonic cue (\texttt{cue\_mode} step or both)}, \\
  \rho_{-k}  &= \textstyle\sum_{t} r_t^{(-k)}
    && \text{total reward on the $k$-th previous trial}, \\
  \alpha_{-k}    &= \1[\text{the $k$-th previous trial contained a decisive lick}], \\
  \varsigma_{-k} &\in \{+1, 0, -1\}
    && \text{rewarded / no response / responded and unrewarded}, \\
  \phi_{-k}      &= \text{first-lick latency on the $k$-th previous trial (s; 0 if none)},
\end{align}
for $k = 1,\dots,K$. $t_{\mathrm{cue}}$ is the step of cue onset (undefined before it, in
which case $c_t = s_t = 0$) and $n_{\mathrm{cue}} = \lceil \texttt{cue\_duration}/\mathrm{d}t
\rceil$. On catch trials every cue channel is zero for the whole trial.
\end{definition}

\begin{remark}[The cue is independent of the trial phase]
$c_t$ depends only on the time since cue onset. If the required delay is shorter than the
cue duration, the cue is still on during the rewarded window.
\end{remark}

\begin{remark}[The previous-trial channels are constant within a trial]
$\rho_{-k},\alpha_{-k},\varsigma_{-k},\phi_{-k}$ are written when a trial ends and held
fixed for every step of the next trial. Cross-trial information therefore reaches the
network through its input, which allows one episode to be one trial.
\end{remark}

\section{Network}\label{sec:net}

The recurrent core is a leaky $\tanh$ network with hidden state $h_t \in \R^{N}$.

\begin{definition}[State equation]
\begin{align}
  a_t &= W_{\mathrm{rec}}\, h_{t-1} + W_{\mathrm{in}}\, u_t + b_{\mathrm{in}}
    + \sqrt{\tfrac{2}{\alpha}}\,\sigma_{\!\xi}\,\xi_t,
    \qquad \xi_t \sim \mathcal{N}(0, I_N), \label{eq:pre}\\[2pt]
  h_t &= (1-\alpha)\, h_{t-1} + \alpha\, \tanh(a_t),
    \qquad h_{-1} = h_0. \label{eq:state}
\end{align}
\end{definition}

\noindent
$\alpha = \mathrm{d}t/\tau$ is the leak, with $\tau$ the unit time constant (a scalar, or
a per-unit vector drawn log-uniformly from a range). $W_{\mathrm{rec}} \in \R^{N\times N}$
is initialised either i.i.d.\ $\mathcal{N}(0, g^2/N)$ or as an orthogonal matrix scaled
by $g$; $(W_{\mathrm{in}})_{jk}\sim\mathcal{N}(0,1/d_u)$ and $b_{\mathrm{in}} = 0$ at
initialisation. The factor $\sqrt{2/\alpha}$ makes the stationary variance of the
noise-driven state independent of $\mathrm{d}t$. Noise is applied in training mode only.
$h_0 = 0$ and is not trained by default.

\begin{definition}[Heads]
\begin{align}
  z_t &= W_{\pi}\, h_t + b_{\pi} \in \R^{2} && \text{policy logits}, \label{eq:logits}\\
  V_t &= w_v^{\top} h_t + b_v \in \R && \text{value estimate}. \label{eq:value}
\end{align}
At initialisation $W_\pi$ is scaled by $0.1$ and $b_\pi = 0$.
\end{definition}

The trained parameters are
$\theta = \{ W_{\mathrm{in}}, b_{\mathrm{in}}, W_{\mathrm{rec}}, W_{\pi}, b_{\pi}, w_v, b_v \}$,
plus $\log\tau$ if \texttt{train\_tau} is set.

\section{Policy}\label{sec:policy}

\begin{definition}[Action distribution]
With $z_t = (z_t^{\mathrm{wait}}, z_t^{\mathrm{lick}})^{\!\top}$,
\begin{equation}
  \pi_t = \softmax(z_t), \qquad A_t \sim \Cat(\pi_t), \qquad A_t = 1 \equiv \text{lick}.
\end{equation}
With a probability floor $\epsilon = \texttt{min\_action\_prob} > 0$ the sampled
distribution is $\tilde\pi_t = \epsilon + (1 - 2\epsilon)\,\pi_t$; the same $\tilde\pi_t$
is used for the log-likelihood below.
\end{definition}

\begin{definition}[Logit gap and lick probability]
\begin{equation}
  \Delta_t = z_t^{\mathrm{lick}} - z_t^{\mathrm{wait}}
    = w_\Delta^{\top} h_t + (b_{\pi}^{\mathrm{lick}} - b_{\pi}^{\mathrm{wait}}),
  \qquad w_\Delta = W_{\pi,\mathrm{lick}} - W_{\pi,\mathrm{wait}},
\end{equation}
\begin{equation}
  p_t = \pi_t(\mathrm{lick}) = \sigma(\Delta_t), \qquad \sigma(x) = (1+e^{-x})^{-1}.
\end{equation}
$p_t$ is a discrete-time hazard rate; the first lick ends the decision, so the lick time
is the first passage of this hazard. $\Delta_t$ is the \emph{policy readout} exported by
\texttt{attach\_states}.
\end{definition}

\begin{definition}[Entropy and log-likelihood]
\begin{equation}
  \mathcal{H}_t = -\!\sum_{a}\! \pi_t(a)\log \pi_t(a), \qquad
  \ell_t = \log \tilde\pi_t(A_t).
\end{equation}
The entropy is logged as a diagnostic and does not enter the loss.
\end{definition}

\section{Environment}\label{sec:env}

Let $\mathcal{P}_t \in \{\textsc{iti}, \textsc{waiting}, \textsc{eligible}, \textsc{post},
\textsc{no\_cue}\}$ be the phase before step $t$ is taken.

\begin{definition}[Accepted lick]
\begin{equation}
  L_t = A_t \cdot \1\!\left[\, t - t_{\mathrm{last}} \ge n_{\mathrm{ref}} \,\right],
  \qquad n_{\mathrm{ref}} = \max\!\bigl(1, \lceil \texttt{lick\_refractory}/\mathrm{d}t \rceil\bigr),
\end{equation}
where $t_{\mathrm{last}}$ is the step of the last accepted lick. Only $L_t$ has
consequences in the environment.
\end{definition}

\begin{definition}[Reward]
\begin{equation}
  r_t = k_{\mathrm{time}}\,\1[\mathcal{P}_t \ne \textsc{post} \;\vee\; \texttt{time\_penalty\_in\_post}]
  \;+\;
  \begin{cases}
    k_{\mathrm{iti}}\, L_t + k_{\mathrm{to}}\,\1[\,n^{\mathrm{iti}}_t \ge n_{\mathrm{to}}\,]
      & \mathcal{P}_t = \textsc{iti},\\[3pt]
    k_{\mathrm{nc}}\, L_t
      & \mathcal{P}_t = \textsc{no\_cue},\\[3pt]
    \gamma_t\, k_{\mathrm{early}}\, L_t
      & \mathcal{P}_t = \textsc{waiting},\\[3pt]
    \gamma_t\, R\, e^{-\kappa \varepsilon_t}\, L_t + k_{\mathrm{miss}}\,\1[\text{window expired}]
      & \mathcal{P}_t = \textsc{eligible},\\[3pt]
    0 & \mathcal{P}_t = \textsc{post}.
  \end{cases}
  \label{eq:reward}
\end{equation}
\end{definition}

\noindent with
\begin{itemize}
  \item $\varepsilon_t = (t - t_{\mathrm{cue}})\,\mathrm{d}t$, the time since cue onset;
  \item $\kappa = \texttt{discount\_rate}$ (0 by default, i.e.\ no temporal discount);
  \item $n^{\mathrm{iti}}_t$ the steps elapsed in the current stop-licking period and
        $n_{\mathrm{to}} = \texttt{iti\_timeout}/\mathrm{d}t$;
  \item $\gamma_t$ the across-trial value gain,
        \begin{equation}
          \gamma_t = \mathrm{clip}\!\left(\exp\bigl(g\,(\rho_{\mathrm{ref}} - \bar\rho_t)\bigr),
                      \gamma_{\min}, \gamma_{\max}\right),
        \end{equation}
        where $\bar\rho_t$ is the rewarded fraction over the last \texttt{reward\_rate\_window}
        completed cued trials of the same environment, $g = \texttt{reward\_rate\_gain}$ and
        $\rho_{\mathrm{ref}} = \texttt{reward\_rate\_ref}$. $g = 0$ gives $\gamma_t = 1$.
        $\gamma_t$ is not observable.
\end{itemize}

\begin{definition}[Phase transitions]
The trial starts in \textsc{iti} with a stop-licking duration drawn from an exponential
distribution of mean \texttt{iti\_mean}, truncated to $[\texttt{iti\_min}, \texttt{iti\_max}]$.
If \texttt{iti\_restart\_on\_lick} is set, an accepted lick redraws the remaining duration.
When the period completes, the trial enters \textsc{waiting} (or \textsc{no\_cue} with
probability \texttt{no\_cue\_prob}) and the cue onset step is recorded. \textsc{waiting}
becomes \textsc{eligible} when $\varepsilon_t$ reaches the required delay. A lick in
\textsc{waiting} or \textsc{eligible} moves the trial to \textsc{post} for
\texttt{post\_lick} seconds, after which the trial ends. The answer window expires
\texttt{answer\_window} seconds after the delay.
\end{definition}

\section{Padding and masks}\label{sec:mask}

Trials in a batch end at different steps. Let $D_t \in \{0,1\}$ be $1$ once the episode
has ended. Two masks are used:
\begin{align}
  m^{\mathrm{data}}_t &= 1 - D_{t-1}, \quad m^{\mathrm{data}}_0 = 1
    && \text{the trial is still running}, \\
  m_t &= m^{\mathrm{data}}_t\,\1[\mathcal{P}_t \ne \textsc{post}]
    && \text{the trial is running and the action has a consequence}.
\end{align}
After termination the observation is frozen and $r_t = 0$. The return
(\S\ref{sec:update}) uses $m^{\mathrm{data}}_t$ so that every reward earned in the trial is
counted; every loss term uses $m_t$ so that the post-lick period, in which no action has
any effect, contributes no gradient. The normaliser for averages over loss-eligible steps is
\begin{equation}
  \mathcal{N} = \sum_{t=0}^{T-1}\sum_{b=1}^{B} m_{t,b}.
\end{equation}

\section{One update}\label{sec:update}

\paragraph{Step 1: roll out.} For $t = 0,\dots,T-1$ and each $b$, run
\eqref{eq:state}, \eqref{eq:logits}, \eqref{eq:value}, sample $A_t$, and step the
environment, storing $\ell_t$, $\mathcal{H}_t$, $V_t$, $h_t$, $r_t$, the masks, and the
indicator $c^{\textsc{iti}}_t = \1[\mathcal{P}_t = \textsc{iti}]$.

\paragraph{Step 2: shape the reward.}
\begin{equation}
  \tilde r_t = c_{\mathrm{scale}}\, r_t - \lambda\, \sg[p_t]\, c^{\textsc{iti}}_t .
  \label{eq:shaped}
\end{equation}
$c_{\mathrm{scale}} = \texttt{reward\_scale}$ rescales the environment reward. $\lambda =
\texttt{iti\_readout\_penalty}$ (0 by default) charges the lick probability at every ITI
step; the stop-gradient keeps this a reward rather than a direct penalty on $p_t$.

\paragraph{Step 3: returns.} With discount $\gamma$, computed backwards from $R_T = 0$:
\begin{equation}
  R_t = \tilde r_t + \gamma\, m^{\mathrm{data}}_t\, R_{t+1}.
\end{equation}

\paragraph{Step 4: advantage.} The value head is the baseline, and the advantage is
standardised over loss-eligible steps:
\begin{equation}
  \hat A_t = \sg\!\left[ R_t - V_t \right], \qquad
  \bar A_t = \frac{\hat A_t - \mu}{s}, \qquad
  \mu = \frac{1}{\mathcal{N}}\sum_{t,b} m\,\hat A, \qquad
  s = \sqrt{\frac{1}{\mathcal{N}}\sum_{t,b} m\,(\hat A - \mu)^2} \vee 10^{-6}.
\end{equation}

\paragraph{Step 5: the loss.}
\begin{align}
  \mathcal{L}_{\pi} &= -\frac{1}{\mathcal{N}} \sum_{t,b} m_t\, \ell_t\, \bar A_t
    && \text{policy gradient (REINFORCE)}, \\
  \mathcal{L}_{V} &= \frac{1}{\mathcal{N}} \sum_{t,b} m_t\, (V_t - R_t)^{2}
    && \text{value regression}, \\
  \mathcal{L}_{a} &= \frac{1}{\mathcal{N}_w} \sum_{t,b} w_t\, \frac{\|h_t\|^2}{N}
    && \text{activity}, \\
  \mathcal{L}_{d} &= \frac{1}{\mathcal{N}_w} \sum_{t,b} w_t\, \frac{\|h_t - h_{t-1}\|^2}{N}
    && \text{activity change}, \\[3pt]
  \mathcal{L} &= \mathcal{L}_{\pi} + c_V\,\mathcal{L}_{V}
                 + \lambda_a\,\mathcal{L}_{a} + \lambda_d\,\mathcal{L}_{d},
\end{align}
where $w_t = m_t\, c^{\textsc{iti}}_t$ when $\texttt{activity\_scope} = \texttt{iti}$ and
$w_t = m_t$ when it is \texttt{full}, $\mathcal{N}_w = \sum_{t,b} w_{t,b}$, $c_V =
\texttt{value\_coef}$, $\lambda_a = \texttt{activity\_penalty}$ and $\lambda_d =
\texttt{activity\_deriv\_penalty}$. The activity terms are differentiable penalties on
the hidden state; they modify the network, not the task.

\paragraph{Step 6: the parameter step.}
\begin{equation}
  g = \nabla_\theta \mathcal{L}, \qquad
  g \leftarrow g \cdot \min\!\left(1, \frac{G}{\|g\|_2}\right), \qquad
  \theta \leftarrow \mathrm{Adam}(\theta, g, \eta),
\end{equation}
with $G = \texttt{grad\_clip}$ ($G = 0$ disables clipping) and $\eta = \texttt{lr}$.

\begin{remark}[Credit assignment]
$\nabla_\theta \mathcal{L}$ is computed by backpropagation through time through the full
recurrence \eqref{eq:state} within the trial. Since $\hat A_t$ is stop-gradiented, the
value head is trained only by $\mathcal{L}_V$, and the policy treats $V_t$ as a constant
baseline.
\end{remark}

\section{Default constants}\label{sec:consts}

\begin{table}[h]\centering
\begin{tabular}{llrl}
\toprule
Symbol & Code name & Default & Meaning \\
\midrule
$\mathrm{d}t$ & \texttt{dt} & $0.02$\,s & step size \\
$\tau$ & \texttt{tau} & $100$\,ms & unit time constant ($\alpha = 0.2$) \\
$N$ & \texttt{hidden} & $128$ & recurrent units \\
$\sigma_{\!\xi}$ & \texttt{noise} & $0.05$ & private noise SD \\
$g$ & \texttt{g} & $1.0$ & recurrent initialisation gain \\
$B$ & \texttt{n\_envs} & $16$ & parallel environments \\
$\eta$ & \texttt{lr} & $4\times10^{-3}$ & Adam learning rate \\
$\gamma$ & \texttt{gamma} & $1.0$ & return discount \\
$c_V$ & \texttt{value\_coef} & $0.5$ & weight on $\mathcal{L}_V$ \\
$G$ & \texttt{grad\_clip} & $1.0$ & gradient-norm clip \\
$c_{\mathrm{scale}}$ & \texttt{reward\_scale} & $0.1$ & reward rescaling \\
$\lambda$ & \texttt{iti\_readout\_penalty} & $0$ & per-step ITI penalty on $p_t$ \\
$\lambda_a$ & \texttt{activity\_penalty} & $0$ ($20$ in \texttt{act}) & activity penalty \\
$\lambda_d$ & \texttt{activity\_deriv\_penalty} & $0$ & activity-change penalty \\
$\epsilon$ & \texttt{min\_action\_prob} & $0$ & action-probability floor \\
\midrule
$R$ & \texttt{reward} & $15.0$ & reward for a lick in the answer window \\
$\kappa$ & \texttt{discount\_rate} & $0$ & temporal discount on the reward \\
$k_{\mathrm{time}}$ & \texttt{time\_penalty} & $-0.05$ & per step \\
$k_{\mathrm{early}}$ & \texttt{early\_penalty} & $-2.0$ & lick during the delay \\
$k_{\mathrm{miss}}$ & \texttt{miss\_penalty} & $-2.0$ & answer window expired \\
$k_{\mathrm{iti}}$ & \texttt{iti\_lick\_penalty} & $-2.0$ & lick in the stop-licking period \\
$k_{\mathrm{nc}}$ & \texttt{no\_cue\_lick\_penalty} & $-2.0$ & lick on a catch trial \\
$k_{\mathrm{to}}$ & \texttt{iti\_timeout\_penalty} & $-10.0$ & stop-licking period timed out \\
$n_{\mathrm{ref}}$ & \texttt{lick\_refractory} & $0.05$\,s & minimum inter-lick interval \\
$n_{\mathrm{cue}}$ & \texttt{cue\_duration} & $0.6$\,s & transient cue duration \\
& \texttt{answer\_window} & $0.8$\,s & answer window \\
& \texttt{post\_lick} & $1.5$\,s & post-lick period \\
& \texttt{iti\_mean}, \texttt{iti\_min}, \texttt{iti\_max} & $3, 2, 5$\,s & stop-licking period \\
& \texttt{no\_cue\_prob} & $0.1$ & catch-trial probability \\
$g$ & \texttt{reward\_rate\_gain} & $1.0$ & across-trial value gain \\
$\rho_{\mathrm{ref}}$ & \texttt{reward\_rate\_ref} & $0.5$ & reference reward rate \\
\bottomrule
\end{tabular}
\caption{Default values from \texttt{variants.BASE}. The complete list, including the
delay-schedule parameters, is generated into \texttt{docs/configuration.md}.}
\end{table}

\end{document}
